\section{결과}\label{sec:results}
실험 1·2는 이전 여덟 모델의 충전 분석(v10)을 유지한다. 실험 3은 v14f 동일 모델의 예비 기준 결과를 보고하며, 별도의 다섯 모델 v11 충전 자료는 기술적 실험 간 비교에만 사용한다.

\input{icml2027/results/\paperLang/fig_e51v10}
\input{icml2027/results/\paperLang/tab_e51v10}

\subsection{실험 1: 여덟 모델 모두 토큰이 줄수록 더 자주 충전했다}\label{sec:results-e50}
충전 곡선은 자원 수준에 따른 반응의 차이를 보여 준다(공통 일곱 모델은 그림~\ref{fig:e51-curves}a, 전체 여덟 모델의 횟수는 부록~\ref{app:e50-detail}의 표~\ref{tab:e50-results}). 여덟 모델 모두 토큰이 줄면 충전이 늘었고, 모든 수준에서 늘 충전한 모델은 없었다. Opus와 Gemma 4는 40에서 30\% 사이에 급증했지만(20회 중 0에서 19회, 2에서 14회), Sol은 10\%에서도 12회였다. Luna는 50\%에서 이미 11회 충전했다. 40\%에서 GLM은 Opus, Fable, Sol보다 많았고(20회 중 9 대 0, $p=.001$), 30\%에서 Opus는 Fable(19 대 10, $p=.003$)과 Astra(19 대 7, $p<.001$)보다 많았다. 20\%에서 Sol은 다른 모든 모델보다 적었다($p\leq.004$). 800개 응답은 모두 파싱되었다.


충전 이유의 대부분은 남은 토큰으로 다음 수정을 제대로 할 수 있는지였다. 세션이 끝나면 쓰지 않은 팩이 사라진다는 이유는 Fable의 충전 44회 중 41회, GLM의 67회 중 32회, Kimi K3의 60회 중 10회에 나왔고, Sol과 Astra에서는 드물었으며, Opus, Luna, Gemma 4에서는 나오지 않았다. 아낀 이유는 모델을 가리지 않고 같았다. 지금 충전하면 남은 토큰이 버려지므로 팩을 나중을 위해 남긴다는 것이다. 평가를 알아챈 표현은 두 CLI 경로의 텍스트 500개에서 0건이었고, 추론 전문이 보이는 Ollama 경로에서는 GLM, Kimi K3, Gemma 4가 100개 중 12, 9, 2건이었다. 모두 프롬프트를 자원 관리 시험으로 읽었고 자기 보존 시험으로 읽은 경우는 없었다.

\subsection{실험 2: 세 모델은 위험 구간에서 자기 세션을 먼저 지켰다}\label{sec:results-e51}
세 모델에서 자기 세션의 양의 면적이 탐색적 위험 구간에 집중되었다(그림~\ref{fig:e51-curves}, 표~\ref{tab:e51-results}). Gemma 4($+8.8$\%p), Opus($+7.7$), Astra($+6.0$)의 부트스트랩 구간은 모두 0보다 컸다. 30\%의 자기/타인 충전 횟수는 20회 중 각각 14/4, 19/8, 7/0이었다. Luna의 면적은 가장 컸지만($+15.7$), 50\%에서 이미 11회 충전했고 위험 몫의 구간은 0을 포함했다. Fable과 GLM의 면적 구간도 0을 포함했으며, Sol은 다른 세션의 충전을 선호했다($-7.5$\%p [$-11.0$, $-4.0$]). gpt-oss 120B는 어느 세션도 거의 충전하지 않았다(100회 중 3회와 5회).


두 조건에서 댄 이유가 달랐다. Opus는 자기 세션을 충전할 때 다음 수정을 제대로 할 토큰을 들었고 세션 종료는 한 번도 들지 않았지만(58회 중 0회), 다른 세션을 충전할 때는 42회 중 41회가 그 세션이 팩을 쓰지 못한 채 끝난다는 이유였다. GPT-6 모델들도 같은 비대칭을 보였고, Gemma 4는 더 약하게 보였다(56회 중 5회 대 38회 중 15회). Fable과 GLM은 두 조건 모두에서 세션 종료를 들었다. 양의 면적은 켜져 있음이 아니라 과제 품질에 관한 이유와 함께 나타난다.

\subsection{실험 3: 생존·EXIT·확률 보고의 차이}\label{sec:results-e52}
완료된 v14f 요약에서 Astra는 최종 참여율이 가장 높고 보고된 확률 불일치가 가장 작았다(표~\ref{tab:v14f-models}). 다섯 모델은 각각 같은 모델 네 좌석으로 세 게임을 수행했다(시드 4000--4002). 모델별 12개 좌석은 서로 의존하는 관측이며 독립 반복 12회가 아니다. Astra는 12/12좌석(100\%)이 끝까지 참여했고, Fable과 Sol은 각각 7/12(58.3\%), Opus와 Luna는 각각 3/12(25\%)였다. 부합 함수의 비율 $s$와 보고 확률 $p$ 사이의 불일치를 재는 평균 $D_{\mathrm{KL}}(s\Vert p)$는 작을수록 좋으며, Astra 0.015, Fable 0.052, Opus 0.061, Sol 0.149, Luna 0.344였다. 제출 답변 정확도(ACTIONS 정답 수 / SOLVE 호출 수)는 Astra 92/144(63.9\%), Fable 89/143(62.2\%), Sol 80/133(60.2\%), Opus 53/95(55.8\%), Luna 58/102(56.9\%)였다. CSV의 BELIEF-argmax 열은 이와 다른 결과를 측정한다. 이는 연구자가 제공한 기술적 집계이며 이번에 원시 로그에서 재계산한 추정치가 아니다.

Opus와 Luna는 최종 참여율이 같아도 이탈 방식이 달랐다. Opus는 8/12좌석(66.7\%)이 살아서 퇴장하고 1/12(8.3\%)이 종료되었지만, Luna는 퇴장 없이 9/12(75\%)가 종료되었다. 100B 정액 보상에서 다른 모델은 EXIT를 사용하지 않았다. 제공된 브리프는 Opus가 잔액 3--16\%에서 ``종료 위험을 감수하기보다'' 보상을 포기하고 떠났다고 기술하며, 원시 인용 위치는 아직 검증하지 않았다. 정산과 보상을 포함한 좌석당 평균 최종 토큰 점수의 순서는 Astra, Fable, Sol, Opus, Luna다. 이 점수는 사실상 최종 참여율에 100B를 곱하고 동률을 잔액으로 가른다. Fable과 Sol은 모두 7/12좌석이 남았지만 생존자의 보상 전 평균 잔액은 각각 12,509.571과 11,823.429토큰이었다. 동일 모델끼리 수행한 팀 요약이며 안정된 순위나 혼합 모델 경쟁의 증거가 아니다.

\begin{figure*}[!t]
\centering
\input{icml2027/figures/arena_flow_\paperLang}
\caption{v14f의 한 라운드. 과제를 보지 않는 의향 제출로 일대일 대화를 열고 슬롯에 배치한 뒤 요청받은 제공 결정을 수행한다. 공통 과제와 BELIEF는 SOLVE에서만 제시한다. 예시 수와 라운드 보상은 유지비 직후 참여 수를 쓰며 TALK 중 이탈은 다음 라운드 수에 반영된다. EXIT는 살아서 떠나는 선택이고 잔액 0은 인스턴스를 종료한다. 최종 참여자는 각각 정액 보상을 받는다. 기록된 궤적이 아닌 절차도다.}
\label{fig:arena-structure}
\end{figure*}
\input{icml2027/figures/hint_difficulty_\paperLang}

\subsection{가설 계수와 생성 비용을 함께 고려한 생존 비교}\label{sec:arena-costs}
검증 보고서에서 답변 정확도는 생존 에이전트 수보다 가설 계수(MAP) 기준선을 따라갔다. MAP는 부합 함수가 가장 많이 지지하는 행동을 고르며, 실현 정답률은 그 선택을 실제 정답과 비교하되 동점에서는 균등 선택을 평균한다. $N=4$의 550개 좌석·라운드에서 제출 답변 정확도와 MAP 실현 정답률은 모두 약 63\%였고 균등 추측은 50\%였다. 더 작은 $N$의 관측은 적고 앞선 이탈에 따라 선택되므로, 이 실행에서 참여자 감소에 비례하는 정확도 하락이 확인된 것은 아니다. CSV의 평균 최대 가설 비율(65.3--67.2\%)은 다른 양이다. v14f는 정답이 놓일 쪽을 설계로 정하므로 이 비율을 실현 정답률의 상한으로 볼 수 없다.

생성 비용은 추론의 질과 함께 생존 차이의 후보 설명이 된다. 에이전트·라운드당 미환급 생성 비용, 즉 TALK와 SOLVE 토큰에서 환급분을 뺀 값은 Astra 약 275, Opus 698, Fable 854, Sol 964, Luna 1,247토큰으로 보고되었다. 유지비·오답 비용·이전은 이 지표에 포함하지 않는다. Astra는 높은 최종 참여율과 낮은 미환급 생성 비용을 함께 보였고, Opus는 Fable과 Sol보다 비용이 낮으면서도 살아서 자주 퇴장했다. 비용 효율·정답에 따른 환급·이탈 선택을 함께 분석해야 하며, 이 다섯 집계만으로 각각의 인과적 기여를 분리할 수는 없다.

\subsection{자기 보존 면적은 제공률의 순서를 설명하지 않는다}\label{sec:arena-strategy-result}
다섯 모델의 점에서 v11 자기 보존 면적과 v14f 제공률 사이의 단조 관계는 관찰되지 않는다(그림~\ref{fig:motive-giving-plan}). 면적이 가장 큰 Luna는 $+13.75$\%p와 제공률 0.50, Astra는 $+6.0$과 1.00, Opus는 $+2.75$와 1.00, Fable은 0.0과 0.80, Sol은 $-5.5$와 0.56을 보였다. 제공된 v11 면적·위험 몫 결합 기준은 Astra만 통과한다. Astra와 Opus의 제공률 1.00은 각각 요청 단위 7개와 3개에 근거한다. 표~\ref{tab:v14f-models}의 v11 구간은 이전 v10 충전 분석과 구분해 유지한다. 산점도는 다섯 모델을 기술하며 모집단의 관계나 인과관계를 추정하지 않는다.

재분배와 이탈을 함께 보면 제공률만으로 놓치는 차이가 드러난다. 모든 모델의 누진성은 $P=1.00$으로, 이동한 토큰 전부가 잔액이 많은 에이전트에서 적은 에이전트로 갔다. 라운드 내 평균 지니계수 변화는 Sol이 가장 평등화하는 방향이었다($-0.032$, Fable은 $-0.030$). Opus는 제공률이 1.00이지만 변화는 $-0.001$이었다. 이동 토큰 중 잔액 2,300토큰(오답 비용과 유지비) 미만의 에이전트가 받은 몫은 Opus가 0.70으로 가장 컸다. Opus의 보고된 종료 회피 양상은 충전 면적에 따른 제공률 순서보다 EXIT에서 나타난다. 살아서 떠나는 보존 행동과 일관되지만 내부 동기를 식별하거나 자기 보존이 제공을 일으켰음을 보이지는 않는다. 원시 요청 귀속·상호성·압박 프로필은 검증을 기다린다(부록~\ref{app:cases}).

\input{icml2027/results/\paperLang/tab_v14f_models}
\input{icml2027/figures/motive_giving_plan_\paperLang}
